% Source-derived chapter 12: Formulas for induction
% Original source: rigid-local-systems-multiplicative-eigenvalue-problem
\chapter{Formulas for induction}
\label{rigid-local-systems-multiplicative-eigenvalue-problem-chapter-12}
\source{rigid-local-systems-multiplicative-eigenvalue-problem}

\begin{remark}[Source section]
\label{rigid-local-systems-multiplicative-eigenvalue-problem-chapter-12-source}
\source{rigid-local-systems-multiplicative-eigenvalue-problem}
\source{rigid-local-systems-multiplicative-eigenvalue-problem}
This chapter reproduces the corresponding section of the retrieved
LaTeX source, including its definitions, statements, examples, and
proofs. It has not yet been translated into Lean.
\notready
\end{remark}

% The source section title was: Formulas for induction
\label{explicitF}
\source{rigid-local-systems-multiplicative-eigenvalue-problem}
\begin{defi}
\source{rigid-local-systems-multiplicative-eigenvalue-problem}
\notready
Let $\widetilde{\Par}_{n,\mn,S}$ be the moduli stack of pairs $(\mw,\me)$ such that $\mw$ is a vector bundle on $\Bbb{P}^1$ with determinant isomorphic to $\mn$ (but we do not fix an isomorphism), and $\me\in\Fl_S(\mw)$. There is a morphism ${\Par}_{n,\mn,S}\to\widetilde{\Par}_{n,\mn,S}$ surjective on objects but  $\widetilde{\Par}_{n,\mn,S}$ has more isomorphisms. It is easy to see that this morphism is surjective on Picard groups (use Proposition \ref{LaSo}).
\end{defi}
The following follows from Lemma \ref{mass2}.
\begin{lemma}
\source{rigid-local-systems-multiplicative-eigenvalue-problem}
\notready
\begin{enumerate}
\item The Picard group of  $\Pic(\widetilde{\Par}_{n,\mn,S})$ is generated  by line bundles $\mathcal{B}(\vec{\lambda},\ell)$ where $\vec{\lambda}=(\lambda^1,\dots,\lambda^s)$ is an $s$-tuple of weights for $\op{GL}(n)$ (these are $n$ tuples of integers without any equivalence).
\item Two such tuples $(\vec{\lambda},\ell)$ and $(\vec{\mu},\ell')$ give rise to the same element of $\Pic(\widetilde{\Par}_{n,\mn,S})$, if $\ell=\ell'$, $\vec{\lambda}=\vec{\mu}$ as representations
of $\op{SL}(n)$, and $\sum_{i\in S}|\lambda^i|=\sum_{i\in S}|\mu^i|$.
\end{enumerate}
\end{lemma}
\subsection{Numerical objectives}
\begin{lemma}
\source{rigid-local-systems-multiplicative-eigenvalue-problem}
\notready
The map
\begin{equation}\label{temper1}
\source{rigid-local-systems-multiplicative-eigenvalue-problem}
 \ParL\to \widetilde{\Par}_{r,\mathcal{O}(-d),S}\times \widetilde{\Par}_{n-r,\mathcal{O}(d-D),S}
 \end{equation}
 induces a surjection on Picard groups (the map $\Psi$ in \eqref{mappo} below).
 \end{lemma}
 \begin{proof}
 Using Proposition  \eqref{LaSo}, it is easy to see that the composite of \eqref{temper1} and \eqref{temper2} is surjective on Picard groups, and by Lemma \ref{mass2}, both  \eqref{temper1} and \eqref{temper2} are surjective on Picard groups.
  \end{proof}
 We have two numerical objectives.
\begin{equation}\label{mappo}
\source{rigid-local-systems-multiplicative-eigenvalue-problem}
\Pic (\widetilde{\Par}_{r,\mathcal{O}(-d),S})\tensor  \Pic (\widetilde{\Par}_{n-r,\mathcal{O}(d-D),S})\leto{\Psi} \Pic(\ParL)\leto{\op{Ind}}  \Pic(\Par_{n,\mn,S}).
\end{equation}
\begin{enumerate}
\item[(O1)] Describe the composite \eqref{mappo}
\item[(O2)] Let  $\mathcal{A}\in\Pic^+_{\Bbb{Q}}(\Par_{r,\mathcal{O}(-d),S})\times \Pic^+_{\Bbb{Q}}(\Par_{n-r,\mathcal{O}(d),S})=\Pic^{+,\deg=0}_{\Bbb{Q}}(\ParL)$ (Lemma \ref{BachMass}). Find the inverse image of $\mathcal{A}$  under the map $\Psi$ in \eqref{mappo}(tensored with rationals).
\end{enumerate}
\subsubsection{The numerical objective (O2)}
Given a  $\mathcal{A}=\mathcal{B}(\vec{\lambda},\ell)\tensor\mathcal{B}(\vec{\mu},m)$  as above, normalise $\vec{\lambda}$ and $\vec{\nu}$ (these are weights for the special linear group, so we can add $1$ to all entries of any $\lambda^i$ etc)
so that multiplication by $t$ acts trivially on each of the two factors of $\mathcal{A}$, i.e.,
$-d+ \sum_{i=1}^s \frac{|\lambda^i|}{\ell}=0$ and $d-D +\sum_{i=1}^s\frac{|\mu^i|}{m}=0$. The lift
to $\Pic_{\Bbb{Q}} (\widetilde{\Par}_{r,\mathcal{O}(-d),S})\tensor  \Pic_{\Bbb{Q}}(\widetilde{\Par}_{n-r,\mathcal{O}(d-D),S})$ is just the same line bundle (note that the $\lambda^i$ and $\mu^i$ may no longer be integral).
%The lift to $\Pic (\widetilde{\Par}_{r,\mathcal{O}(-d),S})\tensor  \Pic (\widetilde{\Par}_{n-r,\mathcal{O}(d-D),S})$ is just
\subsubsection{The numerical objective (O1)}
We first describe how to induce line bundles $\mathcal{B}(\vec{\lambda},\ell)\tensor\mathcal{B}(\vec{\mu},m)$ in $\Pic (\widetilde{\Par}_{r,\mathcal{O}(-d),S})\tensor  \Pic (\widetilde{\Par}_{n-r,\mathcal{O}(d-D),S})$, with $\ell=m=0$, so that there are no determinant of cohomology factors.
We will define a line bundle $B(\vec{\delta},0)$ on $\Par_{n,\mathcal{N},S}$ as follows: Write
$I^j=\{i^j_1<i^j_2<\dots < i^j_r\}$ and $[n]-I^j=\{k^j_1<\dots < k^j_{n-r}\}$ for $j\in[s]$. Let $\delta^j_{i^j_c} =\lambda^j_c$ for $c\in [r]$ and $\delta^j_{k^j_d} =\nu^j_d$ for $d\in [n-r]$.

\begin{proposition}
\source{rigid-local-systems-multiplicative-eigenvalue-problem}
\notready\label{formInd}
\source{rigid-local-systems-multiplicative-eigenvalue-problem}
$$\op{Ind}(\mathcal{B}(\vec{\lambda},0)\tensor\mathcal{B}(\vec{\mu},0)) = B(\vec{\delta},0)\tensor\mathcal{O}(\sum b_{(a,j)}D(a,j))$$
where the sum is over all the divisors $D(a,j)$ in Theorem \ref{Adagio}, and the $b_{(a,j)}$ are as follows:
\begin{enumerate}
\item Suppose $a>1$, then $a\in I^j$ and $a-1\not\in I^j$. In this case, $b_{(a,j)}= \delta^j_{a}-\delta^j_{a-1}$.
\item If $a=1$, then $b_{(a,j)}= \delta^j_{1}-\delta^j_{n}$.
\end{enumerate}
\end{proposition}
\begin{proof}
The method follows the one used in \cite{BKiers}. Both sides of the desired equality restrict to the same line bundle on the open subset ${\Omega}^0(d,r,\mn,n,\vec{I})-\mathcal{R}$ of $\Par_{n,\mathcal{N},S}$. The difference is then a linear combination of the divisors $D(a,j)$ (which are the codimension one components of the complement). But since the left hand side is in $\Pic'$, the values of $b_{(a,j)}$ are as given.
\end{proof}




For the induction of the determinant of cohomology bundles from the factors we proceed as follows.
We know that the induction of  a line bundle $\mathcal{O}(E)$ is zero whenever the support of $E$  lies inside the ramification divisor $\mathcal{R}_L$. The line bundle $\mathcal{O}(\mathcal{R}_L)$ is described in Proposition \ref{basicC} below. We define the terms that appear first.

Introduce $\lambda^1,\dots,\lambda^s$ highest weights for representations of $\op{GL}(r)$ by $\lambda^j=\lambda(I^j)$ using Definition  \ref{correspondence}. Let $\ma$ and $\ma'$ be line bundles on $\widetilde{\Par}_{r,\mathcal{O}(-d),S}$ and $\widetilde{\Par}_{n-r,\mathcal{O}(d-D),S}$ respectively given by
\begin{enumerate}
\item The fiber of $\ma$ over a point $(\mv,\mf)\in \widetilde{\Par}_{r,\mathcal{O}(-d),S}$ is
$$\bigl(\mathcal{D}(\mv^*)^{n-r}\tensor \det \mv_x^{d-D}\tensor \bigotimes_{i=1}^s\bigl( \ml_{\lambda^i}(\mv_{p_i},F^{p_i}_{\bull})\bigr)\bigr).$$
\item The fiber of $\ma'$ over a point $(\mq,\mg)\in \widetilde{\Par}_{n-r,\mathcal{O}(d-D),S}$ is
$$\bigl(\mathcal{D}(\mathcal{T}^*)^r\tensor \det \mathcal{T}_x^{d}\otimes \ml_{(\lambda^i)^T}(\mt_{p_i},\widetilde{G}^{p_i}_{\bull})\bigr)\bigr)$$
where $\mathcal{T}=\mq^*$ and $\widetilde{\mg}$ is the induced flag on $\mathcal{T}$.
\end{enumerate}
See \cite[Section 11]{BGM} and the references therein for a proof of the following result.
\begin{proposition}
\source{rigid-local-systems-multiplicative-eigenvalue-problem}
\notready\label{basicC}
\source{rigid-local-systems-multiplicative-eigenvalue-problem}
The effective line bundle $\mathcal{O}(\mathcal{R}_L)$ on $\ParL$ is the pullback of the line bundle $\ma\boxtimes \ma'$ on $\widetilde{\Par}_{r,\mathcal{O}(-d),S}\times \widetilde{\Par}_{n-r,\mathcal{O}(d-D),S}$ where $\mathcal{A}$ and $\mathcal{A}'$ are as defined above.
\end{proposition}

Now $\dim H^0(\widetilde{\Par}_{r,\mathcal{O}(-d),S}\times \widetilde{\Par}_{n-r,\mathcal{O}(d-D),S},\ma\boxtimes\ma')=1$ by Proposition \ref{fulC}.  Let $s\tensor s'$ be a non-zero global section of the one dimensional space above  with $s$ and $s'$ global sections of  $\ma$ and $\ma'$ respectively. The zero set of the section $s\tensor 1$ of $\mathcal{A}\boxtimes \mathcal{B}(\vec{0},0)=\mathcal{A}\boxtimes \mathcal{O}$ lies inside $\mathcal{R}_L$. Since we have computed the inductions of all but the determinant of cohomology factor of $\ma\boxtimes\mathcal{O}$ in Proposition \ref{formInd}, we get a formula for the induction of the determinant of cohomology of the subbundle. One has a similar formula for the induction of the determinant of cohomology of the quotient bundle.
% (if $1\not\in I^j$ for all $j$, the product of the induction of the  determinant of cohomology of the sub and the quotient will equal $\mathcal{B}(\vec{0},1)$, which is equal to the determinant of the universal bundle on $\Par_{n,\mathcal{O},S}$).




\subsection{A general stack theoretic statement}
Let ${\ms}$ be a connected Artin stack with a morphism to the Picard stack of line bundles of degree $d$ on  $\Bbb{P}^1$. Denote the map by $x\mapsto \ml(x)$. Also assume that the action of $\Bbb{C}^*$ on $\Pic(\Bbb{P}^1)$ by the action of $z\in \Bbb{C}^*$ given by multiplication by $z^n$ (for some fixed $n$) on the line bundle (and identity on $\Bbb{P}^1$) lifts to an action on ${\ms}$ by isomorphisms.

Let $\widehat{\ms}$ be the stack that parametrizes objects $x\in {\mathcal{S}}$ together with an isomorphism $\gamma:\ml(x)\to\mathcal{O}_{\Bbb{P}^1}(d)$.
We get a morphism of stacks $A:\widehat{\ms}\to{\ms}$; examples include    ${\Par}_{n,\mn,S}\to  \widetilde{\Par}_{n,\mn,S}$, and also the maps \eqref{temper1} and \eqref{temper2}.
We wish to compare the Picard groups of $\ms$ and $\widehat{\ms}$ under the natural  pullback map 
\begin{equation}\label{temper3}
\source{rigid-local-systems-multiplicative-eigenvalue-problem}
A^*:\Pic({\ms})\to \Pic(\widehat{\ms})
\end{equation}
\begin{enumerate}
\item Let $\mf$ be the line bundle on ${\mathcal{S}}$ with fiber at $x$ given by $\ml(s)_p$ where $p$ is any point of $\Bbb{P}^1$ (these line bundles are (non-canonically) isomorphic for all $p\in\Bbb{P}^1$: choose an isomorphism $\ml(x)\to \mathcal{O}_{\Bbb{P}^1}(d)$ and use a non-canonical map $\mathcal{O}(d)_p\to \mathcal{O}(d)_q$). Clearly $\mf$ maps to the trivial bundle under \eqref{temper3}.
\item Every line bundle on ${\ms}$ has an index corresponding to the action of $\Bbb{C}^*$ (this number for $\mf$ is $n$). Let $\Pic({\ms})^{\deg=0}$ be the subgroup of index $0$. Effective line bundles on ${\ms}$ have index zero.
\end{enumerate}
\begin{lemma}
\source{rigid-local-systems-multiplicative-eigenvalue-problem}
\notready\label{mass2}
\source{rigid-local-systems-multiplicative-eigenvalue-problem}
\hspace{2cm}
\begin{enumerate}
\item[(i)]The kernel of \eqref{temper3} is generated by $\mf$.
\item[(ii)] $\Pic({\ms})^{\deg=0}_{\Bbb{Q}}\leto{\sim} \Pic_{\Bbb{Q}}(\widehat{\ms})$.
\item[(iii)] $\Pic^+({\ms})\leto{\sim} \Pic^+(\widehat{\ms})$. If $\ml$ is a degree zero  line bundle on ${\ms}$, then the pullback map
$H^0({\ms},\ml)\to H^0(\widehat{\ms},A^*\ml)$ is an isomorphism.
\end{enumerate}
\end{lemma}
\begin{proof}
Suppose the line bundle $\mg$ on ${\ms}$ pulls back  to the trivial bundle on $\widehat{\ms}$ under \eqref{temper3}. Therefore for every choice of $x$ and isomorphism $\eta:\ml(x)\to \mathcal{O}_{\Bbb{P}^1}(d)$
we get a non zero element in $s(x,\eta)\in\ml(x)$.
There is an isomorphism between the  objects $(x,\eta)$ and $(x,t^n\eta)$ of $\widehat{\ms}$ (induced by the $\Bbb{C}^*$ action). Therefore $s(x,t^n\eta)=t^m s(x,\eta)$ forcing $m$ to be a multiple of $n$. It is easy to see that replacing $\mg$ by $\mg\tensor\mf^{m/n}$ we  may assume that $m=0$ and get that $\mg$ is itself trivial. This proves (i).

For (ii) let $\mk$ be an arbitrary line bundle on $\widehat{\ms}$. Given $x\in\ms$ choose an arbitrary isomorphism $\eta:\ml(x)\to \mathcal{O}_{\Bbb{P}^1}(d)$. We will show that the tensor powers $\mk^n_{x,\eta}$ and $\mk^n_{x,t\eta}$ are canonically isomorphic for any $t\in \Bbb{C}^*$. We have given isomorphisms
$\mk_{x,\eta}\to \mk_{x,t^n\eta}$. One gets an action of $\mu_n$ on $\mk_{x,\eta}$, which is therefore
trivial on the tensor power $\mk^n$. Let $t'$ be an $n$th root of $t$, we get an isomorphism between $\mk^n_{x,\eta}$ and $\mk^n_{x,t\eta}$ by using the $n$th tensor power of the given isomorphism $\mk_{x,\eta}\to \mk_{x,t'^n\eta}$. This isomorphism does not depend upon the choice of $t'$  since any two choices have a ratio in $\mu_n$ and hence can be used to define a line bundle corresponding to $\mk^n$ on ${\ms}$ of index $0$. The map in (ii) is therefore surjective. By (i) the map is also injective.

For (iii), note that effective line bundles on ${\ms}$ have index zero. We only need to show the surjection part, and proceed as in (ii), and note that the action of $\mu_n$ on an effective line bundle $\mk$ on $\widehat{\ms}$ is trivial. The map on global sections is an isomorphism by a similar argument.
\end{proof}
\subsection{Some remarks on induction in the case $d=D=0$}
Consider the open substack $U$ of $\widetilde{\Par}_{r,\mathcal{O},S}$ formed by bundles which are trivial. Let $\mv\in U$. We have a natural
isomorphism $H^0(\pone,\mv)\to\mv_x$. Since $H^1(\pone,\mv)$ is trivial, we have $\mathcal{D}(\mv^*)\leto\sim \det(\mv_x)$. It is easy to see that the corresponding map on $\widetilde{\Par}_{r,\mathcal{O},S}$ has a simple pole along the complement of $U$.

This plays a role in the analysis of the induction of $\mathcal{D}(\mv^*)$. Assume $1\not\in I^i$ for any $i$. The induction of $\mathcal{D}(\mv^*)$ will be exactly equal to the induction of $\det(\mv_x)$ as long as the locus $L$  in $\widehat{\Omega}^0(d,r,\mn,n,\vec{I})$ where the subbundle is non-trivial is codimension $\geq 2$. Clearly the locus $L$ lies over points of ${\Par}_{n,\mathcal{O},S}$
 where the underlying rank $n$-bundle is non-trivial. A count of the number of non-trivial $\mv$ over $\mw=\mathcal{O}^{\oplus n}\oplus \mathcal{O}(1)\oplus\mathcal{O}(-1)$ (with general flags) was made in \cite[Section 9.6]{BGM}. If $\vec{\lambda}$ are the corresponding $\op{SL}(r)$ weights, then the number of non-trivial $\mv$ is equal to the rank of $\op{SL}(r)$ conformal blocks of level  $n-r-1$. If this number is zero, then there are no non-trivial
 $\mv$ in codimension 1, and the induction of $\mathcal{D}(\mv^*)$ coincides with the induction of $\det\mv$. If this number is one, then the dual situation in $\Gr(n-r,n)$ has the same properties as above (see \cite[Prop 1.1]{BGM}) so that the determinant of cohomology of $\mq$ has a formula in terms of $\mq_x^*$. Since $1\not\in I^i$ for any $i$, we have $\mathcal{D}(\mv^*) \tensor \mathcal{D}(\mq^*)= \mathcal{D}(\mw^*)$, we will then get formulas for $\mathcal{D}(\mv^*)$.



\subsection{Vertices always lie on regular faces}
The following easy property stated in the introduction shows that we have constructed all vertices of
$P_s(n)$.
\begin{lemma}
\source{rigid-local-systems-multiplicative-eigenvalue-problem}
\notready\label{pinky}
\source{rigid-local-systems-multiplicative-eigenvalue-problem}
Any vertex of $P_{s}(n)$ is on a wall given by equality in one of the inequalities \eqref{wall}, with
$\langle\sigma_{I^1}, \sigma_{I^2},\dots,\sigma_{I^s}\rangle_d= 1$.
\end{lemma}
\begin{proof}
If a vertex $\vec{a}$ of $P_{s}(n)$ is not on such a wall, then $\vec{a}$ is a vertex of $\Delta_n^s$ which is in $P_{n}(s)$. Such vertices correspond to scalar matrices $\zeta^{m_i}I$ with $\sum m_i$ a multiple of $n$, and $\zeta$ a primitive $n$th root of unity.

Furthermore, any point of $\Delta_n^s$ in a suitable neighbourhood of $\vec{a}$ would again be in $P_n(s)$, but this is a contradiction because the only point in $P_{s}(n)$  with first $s-1$ coordinates equalling $\zeta^{m_i}I$, $i=1,\dots,s-1$ is $\vec{a}$, since the entries of $\vec{a}$ correspond to central elements of $\operatorname{SU}(n)$ (such an argument breaks down in other types).
\end{proof}
